% ============================================================================
% LANGENTWURF LE-03a: Las habitaciones (Räume des Hauses)
% Spanisch Klasse 7 - Sequenz 3: Mi casa
% ============================================================================
% Tier: 1 (vollständig)
% Doppelstunde: 1-2 von 8
% Fachfarbe: #c41e3a (Karminrot)
% ============================================================================

\documentclass[11pt,a4paper]{article}

\usepackage[utf8]{inputenc}
\usepackage[T1]{fontenc}
\usepackage[ngerman]{babel}
\usepackage{geometry}
\usepackage{fancyhdr}
\usepackage{lastpage}
\usepackage{xcolor}
\usepackage{tcolorbox}
\usepackage{enumitem}
\usepackage{tabularx}
\usepackage{longtable}
\usepackage{array}
\usepackage{booktabs}
\usepackage{amssymb}

% ============================================================================
% KONFIGURATION
% ============================================================================

\newcommand{\fach}{Spanisch}
\newcommand{\fachfarbe}{c41e3a}
\newcommand{\jahrgangsstufe}{7}
\newcommand{\schulart}{Gesamtschule}
\newcommand{\stundenthema}{Las habitaciones -- Räume des Hauses}
\newcommand{\reihenthema}{Mi casa}
\newcommand{\stundennummer}{3a}
\newcommand{\reihenstunden}{4}
\newcommand{\gerniveau}{A1}
\newcommand{\lehrwerk}{Qué pasa 1}
\newcommand{\lehrwerkseite}{S. 36-39}
\newcommand{\lerngruppengroesse}{24}

% ============================================================================
% LAYOUT
% ============================================================================
\geometry{left=2.5cm, right=2.5cm, top=2.5cm, bottom=2.5cm}
\definecolor{fach}{HTML}{\fachfarbe}
\definecolor{gray}{HTML}{666666}
\definecolor{lightgray}{HTML}{f5f5f5}
\definecolor{successgreen}{HTML}{2a7a4b}
\definecolor{warningorange}{HTML}{b35c00}
\definecolor{basis}{HTML}{2a7a4b}
\definecolor{standard}{HTML}{2c5aa0}
\definecolor{erweiterung}{HTML}{7b1fa2}

\tcbuselibrary{skins,breakable}

\newtcolorbox{infobox}[1][]{
    colback=lightgray,
    colframe=fach,
    fonttitle=\bfseries,
    title=#1,
    breakable
}

\newtcolorbox{scorebox}{
    colback=white,
    colframe=fach,
    fonttitle=\bfseries,
    title=Didaktik-Score,
    breakable
}

\pagestyle{fancy}
\fancyhf{}
\fancyhead[L]{\small\textcolor{gray}{\fach{} -- Jg. \jahrgangsstufe{} -- \stundenthema}}
\fancyhead[R]{\small\textcolor{gray}{LE-\stundennummer{} | Tier 1}}
\fancyfoot[C]{\small\thepage{} / \pageref{LastPage}}
\renewcommand{\headrulewidth}{0.4pt}

% Differenzierungsmarker
\newcommand{\B}{\textcolor{basis}{\textbf{[B]}}}
\newcommand{\Std}{\textcolor{standard}{\textbf{[S]}}}
\newcommand{\E}{\textcolor{erweiterung}{\textbf{[E]}}}

% ============================================================================
% DOKUMENT
% ============================================================================
\begin{document}

% ============================================================================
% DECKBLATT
% ============================================================================
\begin{titlepage}
    \centering
    \vspace*{2cm}
    {\large\textcolor{gray}{\schulart}}\\[2cm]
    {\Huge\textcolor{fach}{\textbf{Langentwurf}}}\\[0.5cm]
    {\large LE-\stundennummer{} | Tier 1 (vollständig)}\\[2cm]
    {\LARGE\textbf{\stundenthema}}\\[0.5cm]
    {\large Sequenz: \reihenthema}\\[2cm]
    \begin{tabular}{rl}
        \textbf{Fach:} & \fach{} (\gerniveau) \\[0.3cm]
        \textbf{Klasse:} & \jahrgangsstufe \\[0.3cm]
        \textbf{Lehrwerk:} & \lehrwerk, \lehrwerkseite \\[0.3cm]
        \textbf{Doppelstunde:} & \stundennummer{} (Std. 1-2 von 8) \\[0.3cm]
    \end{tabular}
    \vfill
\end{titlepage}

\tableofcontents
\newpage

% ============================================================================
% 1. BEDINGUNGSANALYSE
% ============================================================================
\section{Bedingungsanalyse}

\subsection{Lerngruppe}
\begin{tabular}{ll}
    Kursgröße: & \lerngruppengroesse{} SuS \\
    Jahrgangsstufe: & \jahrgangsstufe \\
    Sprachniveau: & \gerniveau{} (GER) -- Anfänger \\
    Fremdsprache: & 2. Fremdsprache (nach Englisch) \\
\end{tabular}

\subsection{Vorwissen aus Sequenz 1-2}
\begin{itemize}
    \item Verb \textit{ser} und \textit{tener} (konjugiert)
    \item Bestimmte Artikel: el, la, los, las
    \item Unbestimmte Artikel: un, una, unos, unas
    \item Pluralbildung: -s / -es
    \item Zahlen 0-20
    \item Redemittel zur Vorstellung und Familie
\end{itemize}

\subsection{Differenzierung (intern)}
\begin{tabular}{|l|p{3.5cm}|p{3.5cm}|p{3.5cm}|}
\hline
& \B{} \textbf{Basis} & \Std{} \textbf{Standard} & \E{} \textbf{Erweiterung} \\
\hline
\textbf{Ziel} & BR: Rezeption + Zuordnung & MR: Produktion mit Hilfe & GY: Freie Beschreibung \\
\hline
\textbf{Vokabeln} & Bild-Wort-Zuordnung & + Artikel einsetzen & + Sätze mit hay \\
\hline
\textbf{Dialog} & Nachsprechen & Mit Stichworten & Frei variieren \\
\hline
\end{tabular}

\vspace{0.5em}
\textbf{WICHTIG:} Diese Marker sind NUR für die Lehrkraft!

% ============================================================================
% 2. SACHANALYSE
% ============================================================================
\section{Sachanalyse}

\subsection{Wortfeld: Las habitaciones}

\textbf{Hauptvokabular (8 Räume):}
\begin{center}
\begin{tabular}{|l|l|l|}
\hline
\textbf{Spanisch} & \textbf{Artikel} & \textbf{Deutsch} \\
\hline
la cocina & la (f.) & die Küche \\
el salón & el (m.) & das Wohnzimmer \\
el dormitorio & el (m.) & das Schlafzimmer \\
el baño & el (m.) & das Badezimmer \\
el comedor & el (m.) & das Esszimmer \\
el jardín & el (m.) & der Garten \\
el garaje & el (m.) & die Garage \\
la terraza & la (f.) & die Terrasse \\
\hline
\end{tabular}
\end{center}

\textbf{Erweiterungsvokabular:}
\begin{itemize}
    \item la habitación (allg. Zimmer)
    \item el pasillo (Flur)
    \item la escalera (Treppe)
    \item el sótano (Keller)
    \item el desván / el ático (Dachboden)
\end{itemize}

\subsection{Grammatische Strukturen}

\textbf{Redemittel für Lokalisierung:}
\begin{itemize}
    \item \textbf{¿Dónde está...?} -- Wo ist...?
    \item \textbf{Está en...} -- Es/Er/Sie ist in...
    \item \textbf{Hay...} -- Es gibt...
\end{itemize}

\textbf{Artikel-Genus:}
\begin{itemize}
    \item Maskulin (-o): el salón, el dormitorio, el baño, el comedor, el jardín, el garaje
    \item Feminin (-a): la cocina, la terraza
    \item Merkhilfe: Endung -a $\rightarrow$ meist \textit{la}, Endung -o/-ón/-or $\rightarrow$ meist \textit{el}
\end{itemize}

\subsection{Kontrastive Betrachtung}

\begin{tabular}{|l|p{5cm}|p{5cm}|}
\hline
& \textbf{Deutsch} & \textbf{Spanisch} \\
\hline
Genus & oft unlogisch (das Mädchen) & meist erkennbar an Endung \\
\hline
``Es gibt'' & mit Singular/Plural & immer \textit{hay} (invariabel) \\
\hline
Wohnzimmer & ein Wort & salón oder cuarto de estar \\
\hline
\end{tabular}

\textbf{Typische Interferenz:} *``el cocina'' (falsches Genus durch dt. ``die'' = la)

\subsection{Kulturelle Aspekte}
\begin{itemize}
    \item Spanische Häuser oft mit Innenhof (\textit{patio})
    \item Terrasse (\textit{terraza}) als wichtiger Lebensraum
    \item Unterschiedliche Wohnformen: piso (Wohnung) vs. casa (Haus)
    \item Siesta-Kultur: Rollos (persianas) in Schlafzimmern
\end{itemize}

% ============================================================================
% 3. DIDAKTISCHE ANALYSE
% ============================================================================
\section{Didaktische Analyse}

\subsection{Stellung in der Sequenz}
\begin{tabular}{|c|l|l|}
\hline
\textbf{Block} & \textbf{Thema} & \textbf{Status} \\
\hline
\textbf{3a} & \textbf{Las habitaciones (Räume)} & \textbf{diese Stunde} \\
3b & Los muebles (Möbel) + hay & folgt \\
3c & Describir mi casa (Beschreibung) & folgt \\
3d & Mi habitación + Sequenztest & folgt \\
\hline
\end{tabular}

\subsection{Kompetenzziele (Rahmenplan MV)}
\begin{itemize}
    \item \textbf{Kommunikativ:} Sprechen (einfache Beschreibungen), Hörverstehen
    \item \textbf{Sprachlich:} Wortschatz (Räume), Grammatik (Artikel)
    \item \textbf{Interkulturell:} Wohnen in Spanien vs. Deutschland
    \item \textbf{Methodisch:} Wortschatzarbeit mit Bildern, Partnerarbeit
\end{itemize}

\subsection{Legitimation}
\textbf{Rahmenplan Spanisch MV (Kl. 7-10):}
\begin{itemize}
    \item Kompetenzbereich: Kommunikative Kompetenz -- Sprechen
    \item Standard: ``Die SuS können ihren Lebensraum mit einfachen Worten beschreiben.''
    \item Themenfeld: ¿Dónde y cómo vivo? (Wohnen)
\end{itemize}

% ============================================================================
% 4. STUNDENZIELE
% ============================================================================
\section{Stundenziele}

\subsection{Grobziel}
Die SuS erweitern ihren \textbf{Wortschatz im Bereich Wohnen}, indem sie die 8 wichtigsten Räume eines Hauses auf Spanisch benennen, die korrekten Artikel verwenden und einfache Lokalisierungsfragen stellen und beantworten.

\subsection{Feinziele (operationalisiert)}
\begin{tabular}{|c|p{7.5cm}|c|c|}
\hline
\textbf{Nr.} & \textbf{Die SuS können...} & \textbf{AFB} & \textbf{Diff.} \\
\hline
FZ1 & die 8 Raumbezeichnungen (cocina, salón, dormitorio, baño, comedor, jardín, garaje, terraza) \textbf{nennen} und korrekt aussprechen & I & alle \\
\hline
FZ2 & den korrekten Artikel (el/la) zu jedem Raum \textbf{zuordnen} und begründen & I & alle \\
\hline
FZ3 & auf die Frage ``¿Dónde está...?'' mit ``Está en...'' \textbf{antworten} & II & alle \\
\hline
FZ4 & einen Hausplan \textbf{beschriften} und kurze Sätze zu Räumen bilden & II & \Std{}/\E{} \\
\hline
FZ5 & einen Mini-Dialog über die eigene Wohnung/das eigene Haus \textbf{entwickeln} & III & \E{} \\
\hline
\end{tabular}

\vspace{1em}
\textbf{AFB-Verteilung:} AFB I: 30\% | AFB II: 45\% | AFB III: 25\%

% ============================================================================
% 5. METHODISCHE ANALYSE
% ============================================================================
\section{Methodische Analyse}

\subsection{Einstieg}
\begin{itemize}
    \item \textbf{Methode:} Bildimpuls -- Foto eines typischen spanischen Hauses
    \item \textbf{Begründung:} Aktivierung von Vorwissen, interkultureller Vergleich
    \item \textbf{Leitfrage:} ``¿Qué ves? -- Was siehst du?''
\end{itemize}

\subsection{Erarbeitung}
\begin{itemize}
    \item \textbf{Methode:} Vokabeleinführung mit Bildkarten + Audio
    \item \textbf{Begründung:} Multimodale Speicherung (visuell + auditiv)
    \item \textbf{Scaffolding:} Farbcodierung für Genus (el = blau, la = rot)
\end{itemize}

\subsection{Übungsphasen}
\begin{itemize}
    \item \textbf{Methode:} Chorisches Sprechen $\rightarrow$ Memory-Spiel $\rightarrow$ Partnerarbeit
    \item \textbf{Progression:} Rezeptiv $\rightarrow$ Reproduktiv $\rightarrow$ Produktiv
    \item \textbf{Differenzierung:} Hausplan mit/ohne Beschriftung
\end{itemize}

\subsection{Medien}
\begin{itemize}
    \item Tafel/Whiteboard: Hausumriss mit leeren Feldern
    \item Lehrwerk: \lehrwerk, \lehrwerkseite
    \item Arbeitsblatt: AB-03a.pdf
    \item Lernpfad: LP-03a.html (Tablets, optional)
    \item Bildkarten: 8 Räume (laminiert)
    \item Audio: Track aus Lehrwerk
\end{itemize}

% ============================================================================
% 6. VERLAUFSPLANUNG
% ============================================================================
\section{Verlaufsplanung}

\textbf{1. Stunde (45 Min.)}

\begin{longtable}{|p{1cm}|p{1.5cm}|p{4cm}|p{3cm}|p{1.5cm}|p{2cm}|}
\hline
\textbf{Zeit} & \textbf{Phase} & \textbf{Lehrerhandeln} & \textbf{Schülerhandeln} & \textbf{SF} & \textbf{Material} \\
\hline
\endhead

0--5 & Einstieg & Begrüßung (span.), Bild eines span. Hauses zeigen: ``¿Qué ves?'' & Vermutungen äußern (dt./span.) & Plenum & Projektor \\
\hline

5--15 & Input & Vokabeln einführen: 8 Räume mit Bildkarten, Audio, Artikel betonen & Hören, Nachsprechen, Notieren & Plenum & Karten, Audio \\
\hline

15--25 & Übung 1 & Chorisches Sprechen: ``El salón -- Alle!'' Genus-Regel erklären & Nachsprechen, Muster erkennen & Plenum & -- \\
\hline

25--35 & Übung 2 & Bild-Wort-Memory in Gruppen anleiten & Spielen, Vokabeln festigen & GA & Memory-Karten \\
\hline

35--45 & Sicherung & Merkkasten auf AB ausfüllen lassen, kontrollieren & AB ergänzen & EA/Plenum & AB-03a \\
\hline
\end{longtable}

\vspace{1em}

\textbf{2. Stunde (45 Min.)}

\begin{longtable}{|p{1cm}|p{1.5cm}|p{4cm}|p{3cm}|p{1.5cm}|p{2cm}|}
\hline
\textbf{Zeit} & \textbf{Phase} & \textbf{Lehrerhandeln} & \textbf{Schülerhandeln} & \textbf{SF} & \textbf{Material} \\
\hline
\endhead

0--5 & Aktivier. & Schnelles Quiz: ``¿Cómo se dice Küche?'' & Antworten: la cocina & Plenum & -- \\
\hline

5--15 & Vertiefung & AB Aufgabe 1-2 anleiten, Zuordnung + Artikel & AB bearbeiten & EA & AB-03a \\
\hline

15--25 & Anwendung & Hausplan beschriften (Aufg. 3), Partnercheck & Beschriften, Vergleichen & EA/PA & AB-03a \\
\hline

25--35 & Dialog & ``¿Dónde está...?'' -- ``Está en...'' einführen, PA üben & Minidialoge führen & PA & -- \\
\hline

35--40 & Erweiter. & \E{}: Kulturvergleich Wohnen in Spanien & Diskutieren, Notizen & Plenum & Tafel \\
\hline

40--45 & Abschluss & Zusammenfassung, HA: Vokabeln lernen, LP-03a optional & Notieren & Plenum & -- \\
\hline
\end{longtable}

\textbf{Legende:} EA = Einzelarbeit, PA = Partnerarbeit, GA = Gruppenarbeit

% ============================================================================
% 7. ERWARTETE SCHWIERIGKEITEN
% ============================================================================
\section{Erwartete Schwierigkeiten}

\begin{tabular}{|p{4.5cm}|p{8cm}|}
\hline
\textbf{Schwierigkeit} & \textbf{Reaktion} \\
\hline
Genus: *el cocina, *la garaje & Farbcodierung (blau/rot), Endungsregel betonen \\
\hline
Aussprache: jardín & Vorsprechen, ``ch'' wie in ``Buch'' (stimmlos) \\
\hline
Verwechslung salón/comedor & Visualisierung: Sofa vs. Esstisch \\
\hline
¿Dónde...? vs. ¿Dónde está...? & Konsequent vollständige Struktur einfordern \\
\hline
Zeitproblem & Puffer: Kulturvergleich kürzen, als HA \\
\hline
\end{tabular}

% ============================================================================
% 8. HAUSAUFGABE
% ============================================================================
\section{Hausaufgabe}

\begin{itemize}
    \item Vokabeln lernen: 8 Räume mit Artikeln
    \item AB-03a fertigstellen (Aufgabe 4: Minidialog)
    \item Optional: Lernpfad LP-03a.html bearbeiten
    \item Optional (\E{}): Eigenes Haus/Wohnung zeichnen und beschriften
\end{itemize}

% ============================================================================
% 9. LÖSUNGEN
% ============================================================================
\section{Lösungen}

\subsection{AB-03a Lösungen}

\textbf{Aufgabe 1: Zuordnung Bild-Wort} (4 Punkte)
\begin{itemize}
    \item Bild Küche $\rightarrow$ la cocina
    \item Bild Wohnzimmer $\rightarrow$ el salón
    \item Bild Schlafzimmer $\rightarrow$ el dormitorio
    \item Bild Badezimmer $\rightarrow$ el baño
\end{itemize}

\textbf{Aufgabe 2: Artikel einsetzen} (4 Punkte)
\begin{enumerate}[label=\alph*)]
    \item \textbf{la} cocina
    \item \textbf{el} jardín
    \item \textbf{la} terraza
    \item \textbf{el} garaje
\end{enumerate}

\textbf{Aufgabe 3: Haus beschriften} (6 Punkte)
\begin{enumerate}
    \item el salón (Wohnzimmer)
    \item la cocina (Küche)
    \item el dormitorio (Schlafzimmer)
    \item el baño (Badezimmer)
    \item el comedor (Esszimmer)
    \item el jardín (Garten)
\end{enumerate}

\textbf{Aufgabe 4: Minidialog} (6 Punkte)

Beispiellösung:
\begin{itemize}
    \item A: ¿Dónde está tu madre?
    \item B: Está en la cocina.
    \item A: ¿Y tu padre?
    \item B: Está en el jardín.
    \item A: ¿Dónde duermes tú?
    \item B: Duermo en el dormitorio.
\end{itemize}

Bewertung: Je 2P pro sinnvollem Austausch (Frage + Antwort)

% ============================================================================
% 10. ANLAGEN
% ============================================================================
\section{Anlagen}

\begin{enumerate}
    \item Arbeitsblatt (AB-03a.pdf)
    \item Musterlösung (ML-03a.pdf)
    \item Lehrerhinweise (LH-03a.pdf)
    \item Lernpfad (LP-03a.html)
    \item Bildkarten: 8 Räume (Kopiervorlage)
    \item Memory-Karten (Kopiervorlage)
    \item Hausumriss (Tafelbild)
\end{enumerate}

% ============================================================================
% 11. DIDAKTIK-SCORE
% ============================================================================
\newpage
\section{Didaktik-Score}

\begin{scorebox}
\textbf{Bewertung der didaktischen Qualität nach 10 Dimensionen}\\
\textbf{Ziel-Score: $\geq$ 9.0 (Premium-Qualität)}

\vspace{1em}
\begin{tabular}{|c|l|c|c|p{4cm}|}
\hline
\textbf{\#} & \textbf{Dimension} & \textbf{Gew.} & \textbf{Score} & \textbf{Notiz} \\
\hline
1 & Differenzierung & 15\% & 9/10 & [B]/[S]/[E] qualitativ \\
\hline
2 & Taxonomie (AFB) & 10\% & 9/10 & 30/45/25 erfüllt \\
\hline
3 & Lernziel-Alignment & 10\% & 10/10 & RP MV explizit \\
\hline
4 & Aktivierung & 10\% & 9/10 & Memory, Dialog \\
\hline
5 & Scaffolding & 10\% & 9/10 & Farbcodierung, gestufte Hilfen \\
\hline
6 & Feedback-Qualität & 10\% & 9/10 & LP: elaboriert \\
\hline
7 & Sprachliche Klarheit & 10\% & 10/10 & Kl. 7 angemessen \\
\hline
8 & Motivation & 10\% & 9/10 & Alltagsbezug Wohnen \\
\hline
9 & Inklusion (UDL) & 10\% & 9/10 & Visuell + auditiv \\
\hline
10 & Praktikabilität & 5\% & 9/10 & 90 Min. realistisch \\
\hline
\multicolumn{3}{|r|}{\textbf{GESAMT:}} & \textbf{9.1/10} & \\
\hline
\end{tabular}

\vspace{1em}
$\checkmark$ \textcolor{successgreen}{\textbf{FREIGABE} (Score $\geq$ 9.0)}
\end{scorebox}

\end{document}
